\documentclass{article}

\begin{document}
Planificación del proyecto organizada en sprints, con tareas específicas para cada uno de ellos.
Sprins de dos semanas de duración, reuniones cada dos semanas revisando el progreso y ajustando 
la planificación según sea necesario.
    
\begin{itemize}
    \item \textbf{Sprint 1: Investigación inicial (21/10/2025)}
    \begin{itemize}
        \item Investigación sobre PyTorch y tensores para procesado de imágenes
        \item Investigación sobre imágenes en formato crudo y estructura interna
        \item Investigación sobre procesado de imágenes
        \begin{itemize}
            \item Espacios de color (RGB, HSV)
            \item Debayering
        \end{itemize}
        \item Investigación sobre la librería de interfaz gráfica PyQt6
        \item Preparación del entorno de desarrollo:
        \begin{itemize}
            \item Configuración de repositorio Git
            \item Configuración de entorno virtual con dependencias necesarias
            \item Configuración de IDE de desarrollo (VSCode)
            \item Instalación de CUDA y PyTorch para desarrollo con GPU
        \end{itemize}
    \end{itemize}
    \item \textbf{Sprint 2: Primeras implementaciones (04/11/2025)}
    \begin{itemize}
        \item Investigación sobre funcionamiento de programas de procesado de imágenes
        similares
        \item Carga y tratamiento de imágenes con tensores de PyTorch
        \item Debayering de imágenes RAW
        \item Ajuste de contraste mediante función sigmoide
        \item Visualización de resultados con matplotlib
    \end{itemize}
    \item \textbf{Sprint 3: Esqueleto inicial (18/11/2025)}
    \begin{itemize}
        \item Interfaz básica con PyQt6
        \item Integración de carga de imágenes en la interfaz
        \item Visualización de imágenes en la interfaz
        \item Diseño de arquitectura
        \begin{itemize}
            \item Definición de módulos y clases principales
            \item Core
            \item GUI
            \item Loaders
        \end{itemize}
    \end{itemize}
\end{itemize}

\newpage 

\begin{itemize}
    \item \textbf{Sprint 4: Implementaciones avanzadas I (2/12/2025)}
    \begin{itemize}
        \item Diseño para la definición de \textit{pipelines} de procesado
        \item Gestión completa de imagen en la aplicación
        \begin{itemize}
            \item Carga
            \item Visualización
            \item Aplicación de operaciones de procesado
            \item Reseteo a estado original
        \end{itemize}
        \item Implementación de operaciones de procesado I
        \item Implementación de histograma RGB
        \item Exportación de \textit{pipelines} de procesado
    \end{itemize}
\end{itemize}

\begin{itemize}
    \item \textbf{Sprint 5: Implementaciones avanzadas II (16/12/2025)}
    \begin{itemize}
        \item Implementación de operaciones de procesado II
        \item Diseño e implementación de sistema de \textit{snapshots} para mejorar
        el rendimiento durante el procesado de la imagen 
        \item Optimización de memoria durante el tratamiento de la imagen
        \item Tests unitarios para las operaciones de procesado implementadas 
        \item Revisión de gestión y manejo de errores en la aplicación
        \item Implementación de módulo de \textit{benchmarking}
    \end{itemize}
    \item \textbf{Sprint 6: Implementaciones avanzadas III (13/01/2026)}
    \begin{itemize}
        \item Implementación de operaciones de procesado III
        \item Implementación de carga de imágenes en formato propio (PKL.GZ)
        \item Implementación de gráfico de perfiles de intensidad
        \item Guardado de imagen en formato JPG/JPEG
        \item Selección de dispositivo de procesamiento (CPU/GPU)
        \item Implementación de \textit{logger} para seguimiento de eventos
    \end{itemize}
    \item \textbf{Sprint 7: Revisión y refinamiento (27/01/2026)}
    \begin{itemize}
        \item Revisión de interfaz gráfica y experiencia de usuario
        \begin{itemize}
            \item Mejoras en visualización y navegación del visor de imágenes
            \item Mejoras y refinamiento de controles para ajuste de parámetros de procesado
        \end{itemize}
        \item Implementación y refinamiento de tests unitarios para nuevas operaciones
        \item Ajuste y refinamiento de sistema de zoom
        \item Revisión y refactorización de código
    \end{itemize}
\end{itemize}



\end{document}